\documentclass[10pt,a4paper]{amsart}
\usepackage[margin=19mm]{geometry}
\usepackage{amsmath,amssymb,mathtools}
\usepackage[scr=rsfs,cal=euler]{mathalfa}
\usepackage{xfrac}
\usepackage[unicode,hidelinks,bookmarks=false]{hyperref}
\newcommand{\ie}{i.e.\ }
\newcommand{\cf}{cf.\ }
\newcommand{\resp}{respectively\ }
\newcommand{\Subsection}{\subsection}
\providecommand{\nobreakdash}{\nobreak\hbox{-}}
\setlength{\emergencystretch}{2em}
\multlinegap0pt
\usepackage{bm}
\usepackage{bbm}
\usepackage{dsfont}
\usepackage{xspace}
\usepackage{cancel}
\usepackage{mathtools}
\usepackage{amsthm}
\makeatletter
\@ifundefined{thm}{\newtheorem{thm}{Thm}}{}
\@ifundefined{lem}{\newtheorem{lem}[thm]{Lemma}}{}
\makeatother
\title[A Reshetnyak type theorem for quasiregular values on Carnot ]{A Reshetnyak type theorem for quasiregular values on Carnot groups of $H$-type}
\author{Deguang Zhong}
\date{Source: arXiv:2509.21025v1 (CC BY 4.0)}
\begin{document}
\maketitle

\noindent\emph{Source and licence.} A Reshetnyak type theorem for quasiregular values on Carnot groups of $H$-type. Version arXiv:2509.21025v1 (CC BY 4.0), distributed under the Creative Commons Attribution 4.0 International licence, as recorded in the arXiv licence metadata for this exact version and shown on the abstract page https://arxiv.org/abs/2509.21025v1. Full rights notice: licenses/arxiv-2509.21025-CC-BY-4.0.txt.\par\smallskip

\noindent\textit{Editorial scope.} Counted unit: the Lem whose source label is \texttt{lemma3522} in the article (source0.tex lines 971--976, source label \texttt{lemma3522}), delivered together with the article's own complete original proof of it (source0.tex lines 977--1007, one \texttt{proof} environment of 31 lines). No mathematics is written, repaired or re-derived: statement and proof are the article's own text line for line. Required context retained uncounted: -11 (lines 551--557); 0311888888000 (lines 553--555); 101209 (lines 852--856). Every definition, display and label of those lines is retained unchanged. Omitted from this package and not redistributed: 1--550, 556--851, 857--970, 1008--1621. Nothing inside a retained excerpt is omitted or altered: every retained body is delivered exactly as the article's lines are written, so no omission receipt of any kind accompanies this package. Citations: the retained excerpts carry no citation command at all, so no bibliography entry of the article is transcribed and no reference is redistributed. Reference adaptation: none was needed and none was made; every reference inside the delivered unit resolves inside it, so no unresolved reference or citation is produced. Printed slips kept literal: no printed word, symbol, hypothesis or number of the counted statement or of its proof was altered. Layout and class: the article's own class is \texttt{article} with packages including amsfonts, amssymb, amsmath, balance, amsthm, url, hyperref; the wrapper is the lane's pinned \texttt{amsart} editorial wrapper at 10pt on A4, which restates the theorem-like environments the retained text needs and adds the article's own macro definitions that the retained text needs (none), with extra wrapper packages (bm, bbm, dsfont, xspace, cancel, mathtools). The delivered numbering is the unit's own. Assets: no retained span contains a figure environment, a graphics-inclusion command, a PDF-page inclusion, a table or a TikZ picture; no image file is copied into this package, the package declares no asset dependency and no image file is redistributed with it. Licence: version arXiv:2509.21025v1 of this article is distributed under the Creative Commons Attribution 4.0 International licence, as recorded in the arXiv licence metadata for this exact version and shown on the abstract page https://arxiv.org/abs/2509.21025v1. A full rights notice is delivered at licenses/arxiv-2509.21025-CC-BY-4.0.txt. Characters: every retained excerpt is pure ASCII, so the delivered engine, XeLaTeX with its default Latin Modern text font, reports no missing character.
\begin{lem}\label{-11}
Suppose that $\Omega\subset\mathbb{G}$ is a bounded domain and $p\geq1$. If $f\in W^{1,p}(\Omega),$ then there exists a constant $C>0,$ such that the following inequality
\begin{equation}\label{0311888888000}
\left(\int_{B(x,r)}\vert f(x)-f_{B(x,r)}\vert^{p}dx\right)^{1/ p}\leq Cr\left(\int_{B(x,r)}\vert \nabla_{h}f(x)\vert^{p}dx\right)^{1/p}
\end{equation}
holds for all $x\in\Omega$ and $r>0$ such that $B(x,r)\subset\Omega.$
\end{lem}

\begin{lem}\label{101209}
Let $\varepsilon_{0}\in(0,1)$ be a given constant. Suppose that $\mathbb{G}$ is a Carnot group of $H$-type,  $y_{0}\in\mathbb{G},$ $\Omega\subset\mathbb{G}$ is a domain and $f:\Omega\rightarrow \mathbb{G}$ is a mapping of class $W_{loc}^{1,Q}(\Omega)$ satisfying  the following inequality
$$\vert D_{h}f(x)\vert^{Q}\leq K(x)J_{f}(x)+\Sigma(x)(\rho\circ l_{y_{0}^{-1}}\circ f(x))^{Q}$$
for almost every $x\in\Omega.$ Let $K\in L_{loc}^{p}(\Omega)$ and $\Sigma\in L_{loc}^{1+\varepsilon}(\Omega),$ where $(1-\varepsilon_{0})(1+\varepsilon)/((1-\varepsilon_{0})\varepsilon-\varepsilon_{0})<p\leq\infty$ and $\varepsilon_{0}/(1-\varepsilon_{0})<\varepsilon\leq\infty.$  If we further assume that $K$  satisfies $(\mathcal{K},Q-1)$-spherical  condition on $\Omega,$ then either $f\equiv y_{0},$ or $\mathcal{H}^{1}(f^{-1}\{y_{0}\})=0.$  In the latter case, $f^{-1}\{y_{0}\}$ is totally disconnected.
\end{lem}

\begin{lem} \label{lemma3522}
For a given $\varepsilon_{0}\in(0,1),$ we suppose that $\varepsilon_{0}/(1-\varepsilon_{0})<\varepsilon\leq\infty$ and $(1-\varepsilon_{0})(1+\varepsilon)/((1-\varepsilon_{0})\varepsilon-\varepsilon_{0})<p\leq\infty.$ Let $\mathbb{G}$ be a Carnot group of $H$-type,  $y_{0}\in\mathbb{G},$ $\Omega\subset\mathbb{G}$ be a domain and $f:\Omega\rightarrow \mathbb{G}$ be an nonconstant mapping of class $W_{loc}^{1,Q}(\Omega)$ satisfying  the following inequality
$$\vert D_{h}f(x)\vert^{Q}\leq K(x)J_{f}(x)+\Sigma(x)(\rho\circ l_{y_{0}^{-1}}\circ f(x))^{Q}$$
for almost every $x\in\Omega.$ Suppose that $K\in L_{loc}^{p}(\Omega),$  $\Sigma\in L_{loc}^{1+\varepsilon}(\Omega)$  and $K$ satisfies $(\mathcal{K},Q-1)$-spherical  condition on $\Omega.$ If $\vert D_{h}f\vert^{Q}/K(\rho\circ l_{y_{0}^{-1}}\circ f)^{Q}\in L_{loc}^{1}(\Omega),$
then $\log\rho\circ l_{y_{0}^{-1}}\circ f\in W_{loc}^{1,q_{1}}(\Omega),$ where $q_{1}=pQ/(1+p)<Q.$ In particular, $\log\rho\circ l_{y_{0}^{-1}}\circ f\in W_{loc}^{1,Q}(\Omega)$ if $K(x)\equiv K\geq1.$
\end{lem}

\begin{proof}
Without loss of generality, we here only consider the case for $y_{0}=0,$ $K\in L^{p}(\Omega),$  $\Sigma\in L^{1+\varepsilon}(\Omega)$ and $\vert D_{h}f\vert^{Q}/K(\rho\circ l_{y_{0}^{-1}}\circ f)^{Q}\in L^{1}(\Omega).$ By letting $g_{\lambda}:=\log\vert f\vert_{\lambda}:\Omega\rightarrow\mathbb{R}$ for every $\lambda>0,$ we see that $g_{\lambda}\in W^{1,Q}(\Omega).$ For every fixed $x_{0}\in\Omega,$ we chose $B=B(x_{0},r)$ which compactly contained in $\Omega.$ Since 
\begin{equation}\label{1009982887}
\int_{B}\vert\nabla_{h}g_{\lambda}\vert^{q_{1}}\leq\int_{B}\frac{\vert D_{h}f\vert^{q_{1}}}{(\rho\circ f)^{q_{1}}}=\int_{B}\left(\frac{\vert D_{h}f\vert^{Q}}{K(\rho\circ f)^{Q}}\right)^{\frac{q_{1}}{Q}}K^{\frac{q_{1}}{Q}}\leq C,
\end{equation}
where $C=C(Q,q_{1},\Omega,\mathcal{K}, \| \vert D_{h}f\vert^{Q}/K(\rho\circ l_{y_{0}^{-1}}\circ f)^{Q}\|_{L^{1}(\Omega)})$ is a constant does not depend on $\lambda,$ we see by Poincar\'{e} inequality in Lemma \ref{-11}  that 
 $$\int_{B}\vert g_{\lambda}-(g_{\lambda})_{B}\vert dx
\leq C(B)\int_{B}\vert g_{\lambda}-(g_{\lambda})_{B}\vert^{q_{1}} dx\leq C(B)\int_{B}\vert \nabla_{h}g_{\lambda}\vert^{q_{1}} dx\leq C,$$
 where $C=C(Q,q_{1},\Omega,\mathcal{K}, \| \vert D_{h}f\vert^{Q}/K(\rho\circ l_{y_{0}^{-1}}\circ f)^{Q}\|_{L^{1}(\Omega)})$ is a constant does not depend on $\lambda.$ 
 Without loss of generality, we may assume that $\rho\circ f(x)<1$ and $0<\lambda<1.$ In addition, from the proof of Lemma  \ref{101209},  we see that $\log\log\frac{1}{\rho}\circ f\in W_{loc}^{1,q_{1}}(\Omega),$ where $q_{1}=pQ/(1+p).$ Hence, $f^{-1}\{y_{0}\}$ has zero Haar measure. So we may assume that $\rho\circ f(x)$ is not identically $0$ on $B.$ Then, there exists $0<t_{0}<t_{1}<1,$ such that the set $A:=\{x\in B:t_{0}<\rho\circ f(x)<t_{1}<1\}$ has positive  measure. In addition, 
 \begin{equation}
\begin{aligned}
&\int_{B}\vert g_{\lambda}-(g_{\lambda})_{B}\vert dx\geq\int_{A}\vert g_{\lambda}-(g_{\lambda})_{B}\vert dx\\
\geq&\int_{A}(\vert (g_{\lambda})_{B}\vert-\vert g_{\lambda}\vert )dx\geq\int_{A}(\vert (g_{\lambda})_{B}\vert-\log\frac{1}{t_{0}})dx.
\end{aligned}
\end{equation}
  Hence,
  \begin{equation}\label{0009j}
  \int_{B}\vert g_{\lambda}\vert dx\leq C(m(A),B,t_{0},Q).
  \end{equation}
By monotone convergence theorem, we get that $\log\rho\circ f\in L^{1}(B).$ We next will show that $\log\rho\circ f\in  L^{q_{1}}(B).$ By Poincar\'{e} inequality   (\ref{0311888888000})  and inequalities (\ref{1009982887}) and (\ref{0009j}), we get 
\begin{equation}
\begin{aligned}
&\int_{B}\vert g_{\lambda}\vert^{q_{1}}dx\\
\leq& C(q_{1})\int_{B}\vert g_{\lambda}-(g_{\lambda})_{B}\vert^{q_{1}}dx+C(q_{1})\int_{B}\vert(g_{\lambda})_{B}\vert^{q_{1}}dx\\
\leq& C(B,q_{1})\int_{B}\vert \nabla_{h}g_{\lambda}\vert^{q_{1}}dx+C(B,q_{1})\left(\int_{B} \vert\log\rho\circ f\vert dx\right)^{q_{1}}\leq C,\\
\end{aligned}
\end{equation}
where $C$ depends only on $B,q_{1},Q $ and $\| \log\rho\circ f\|_{L^{1}(B)}.$ Hence, by  monotone convergence theorem again, we get that $\log\rho\circ f\in L^{q_{1}}(B).$ In addition,  monotone convergence Theorem  and inequality (\ref{1009982887}) make sure that 
$\nabla_{h}\log\rho\circ f\in L^{q_{1}}(B).$ Therefore, we have $\log\log\frac{1}{\rho}\circ f\in W^{1,q_{1}}(B),$ where $q_{1}=pQ/(1+p).$ This finishes the proof. 
\end{proof}

\end{document}
